
\scanpage{56}
\mspage{12}
\deleted{\unsure{A)}} \inserted{(B)}
\uline{Opération} de \deleted{\uline{\unsure{raffinement} de}} \inserted{sur}\uline{marquage}

\deleted{On peut le définir \inserted{en termes de} par la donnée, d'une \unsure{composante connexe de} pour chaque arête du contour $R$, \unsure{$R\ne\varnothing$} $\Sigma\setminus R$. Soit \unsure{$\widetilde\Sigma_*$} $A_0^+\in\unsure{A_0}$ l'ens}
\deleted{\unsure{des}}
Soit $A^+$ l'ens des arêtes \inserted{(libres ou liées)} ``effectivement marquées'', \inserted{i.e. telles que $r(a)>0$},
\deleted{et $B$ l'ens des},
correspondant à des $C\in\pi_0(\Sigma)$ telles que $C\cap R\ne\varnothing$, soit $B^+$ l'ens des comp. connexes de
\[
 \Sigma^+\setminus R\cap\Sigma^+
\]
(lesquelles sont des int. ouverts), un surmarquage est défini par la donnée d'une application
\[
 B^+\xrightarrow{r'}\mathbb N,
\]
et peut se décrire \deleted{par} \inserted{comme} l'opération qui sur chaque
\uline{arête} \inserted{\unsure{(ouverte)} \deleted{\unsure{qui}}}
de $\Sigma$, correspondant à un $b\in B^+$, $r'(b)$ nouveaux sommets de \deleted{\unsure{marquage}} \inserted{rigidification}.
Cette fois-ci, \deleted{\unsure{l'opé}} il y a des \unsure{choix} arbitraires --- néanmoins l'opération donne un objet de $S\mathcal TD^+$, défini à isom unique près \inserted{soit $(X,\Sigma,R_1)$} --- muni bien entendu d'un ens marqué $R'\subset R_1$, tel que l'opération inverse de gommage de $R'$ redonne l'objet de départ $(X,\Sigma,R)$.

\margin{\deleted{arêtes} \deleted{\unsure{effectivement}} cercles marqués, \unsure{comme} \uline{arêtes} sur $\Sigma$}
\margin{(\uline{NB} \unsure{elle est réunion|elles se déduisent} des cercles eff. marqués, et non des intervalles ouverts)}
\mspage{55}
